\begin{example}[A field]
Let $K$ be a field.  Since $\Spec K=\{(0)\}$, for $a\in K$,
\[
D(a)=\begin{cases}
\varnothing,&a=0,\\
\Spec K,&a\neq0.
\end{cases}
\]
This is the unit/nilpotent dichotomy in its simplest possible form.
\end{example}

\begin{example}[The punctured affine line]
In $A=k[t]$,
\[
D(t)=\Spec k[t]\setminus V(t).
\]
If $k$ is algebraically closed, this removes the closed point $(t)$ and keeps the generic
point together with all $(t-a)$ for $a\neq0$.  Algebraically,
\[
A_t=k[t,t^{-1}],
\qquad
D(t)\cong\Spec k[t,t^{-1}].
\]
\end{example}

\begin{example}[Removing two points from an affine line]
Assume $a\neq b$ in an algebraically closed field $k$.  Then
\[
D((t-a)(t-b))=D(t-a)\cap D(t-b)
\]
is the affine line with the two closed points $(t-a)$ and $(t-b)$ removed.  The coordinate
ring of this principal open is
\[
k[t]_{(t-a)(t-b)}.
\]
\end{example}

\begin{example}[Arithmetic localization]
In $\Spec\ZZ$,
\[
D(6)=\{(0)\}\cup\{(p):p\neq2,3\}.
\]
The points $(2)$ and $(3)$ are removed because they contain $6$.  Correspondingly,
\[
D(6)\cong\Spec\ZZ[1/6].
\]
The localization makes both $2$ and $3$ units.
\end{example}

\begin{example}[Coordinate axes removed]
In $k[x,y]$,
\[
D(xy)=D(x)\cap D(y).
\]
Thus $D(xy)$ is the region where both coordinates are invertible, and
\[
k[x,y]_{xy}\cong k[x,x^{-1},y,y^{-1}].
\]
Geometrically it is the affine plane with both coordinate axes removed.
\end{example}

\begin{example}[A principal open need not behave additively]
In $k[x,y]$,
\[
D(x+y)
\]
is the complement of the diagonal line $x+y=0$.  It is not equal to
$D(x)\cup D(y)$: the prime $(x,y)$ lies in neither side, while the prime $(x+y)$ lies in
$D(x)\cup D(y)$ but not in $D(x+y)$.  This illustrates why multiplication, not addition,
is the natural operation for principal opens.
\end{example}

\begin{example}[A two-element principal-open cover]
In $k[t]$,
\[
(t)+(1-t)=k[t],
\]
so
\[
D(t)\cup D(1-t)=\Spec k[t].
\]
Indeed no prime ideal can contain both $t$ and $1-t$, because then it would contain $1$.
The identity
\[
t+(1-t)=1
\]
is the algebraic certificate of the cover.
\end{example}

\begin{example}[An open set that is not principal]
In $A=k[x,y]$,
\[
U=D(x)\cup D(y)=\Spec A\setminus V(x,y).
\]
This is the complement of the origin.  It is a union of two principal opens, but in general
it is not itself principal.  If $U=D(f)$, then $V(f)=V(x,y)$, hence
$\sqrt{(f)}=(x,y)$; over a polynomial ring in two variables this would make a height-two
maximal ideal the radical of a principal ideal, contradicting the principal ideal theorem.
Thus a basis element is not closed under arbitrary unions.
\end{example}

\begin{example}[Crossing axes: localization removes a component]
Let
\[
R=k[x,y]/(xy).
\]
Then on $D(\bar x)$, the element $\bar x$ is invertible and
$\bar x\bar y=0$ forces $\bar y=0$.  Therefore
\[
R_{\bar x}\cong k[x,x^{-1}].
\]
Similarly $R_{\bar y}\cong k[y,y^{-1}]$.  Localization discards the component on which the
inverted function vanishes identically.
\end{example}

\begin{example}[Crossing axes minus the origin]
Again let $R=k[x,y]/(xy)$ and localize at $\bar x+\bar y$.  In $R_{\bar x+\bar y}$ set
\[
e=\frac{\bar x}{\bar x+\bar y},
\qquad
1-e=\frac{\bar y}{\bar x+\bar y}.
\]
Since $\bar x\bar y=0$, one checks $e^2=e$.  The two idempotent pieces separate the two
punctured axes, giving
\[
R_{\bar x+\bar y}\cong k[x,x^{-1}]\times k[y,y^{-1}].
\]
Thus removing the intersection point disconnects the reducible curve.
\end{example}

\begin{example}[A quotient and a principal open]
Let $A=k[x,y]$, $I=(xy)$, and take $f=x+y$.  Then
\[
D_{A/I}(\overline{x+y})
\cong
V_A(xy)\cap D_A(x+y).
\]
The left side is computed by
\[
(A/I)_{\overline{x+y}}\cong A_{x+y}/(xy)A_{x+y}.
\]
This is the union of the two axes with their intersection point removed.
\end{example}

\begin{example}[Preimage of a principal open]
Consider $\varphi:k[t]\to k[x]$ with $t\mapsto x^2$.  The induced map
\[
\varphi^\ast:\Spec k[x]\to\Spec k[t]
\]
satisfies
\[
(\varphi^\ast)^{-1}(D(t))=D(x^2)=D(x).
\]
Thus the inverse image of the punctured $t$-line is the punctured $x$-line.
\end{example}

\begin{example}[Product rings]
Let $R=A\times B$ and $(a,b)\in R$.  Under
\[
\Spec(A\times B)\cong\Spec A\sqcup\Spec B,
\]
one has
\[
D_{A\times B}(a,b)
\cong
D_A(a)\sqcup D_B(b).
\]
For the idempotent $(1,0)$, this gives
\[
D(1,0)\cong\Spec A,
\]
while the entire $\Spec B$ component is removed.
\end{example}

\begin{example}[Dual numbers]
Let
\[
A=k[\varepsilon]/(\varepsilon^2).
\]
The element $\varepsilon$ is nilpotent, so
\[
D(\varepsilon)=\varnothing.
\]
But $1+\varepsilon$ is a unit with inverse $1-\varepsilon$, so
\[
D(1+\varepsilon)=\Spec A.
\]
The one-point topology sees the two extreme principal opens but does not itself record the
nilpotent thickness.
\end{example}
